\documentclass[10pt]{extarticle}

\usepackage{custom}
\usepackage[sfdefault]{atkinson}
\usepackage[T1]{fontenc}

\title{Project List 1}
\author{33-467: Astrophysics of the Stars and Galaxy}
\date{Fall 2023}

\begin{document}
\maketitle

\definecolor{tcol_CNT1}{HTML}{EDEADE} % First color for Contents
\definecolor{tcol_CNT2}{HTML}{EDEADE} % Second color for Contents
\definecolor{tcol_CNV1}{HTML}{8E44AD} % First color for Conventions
\definecolor{tcol_CNV2}{HTML}{A10B49} % First color for Conventions

\begin{tcolorbox}[enhanced,
    fonttitle=\fontsize{20}{24}\sffamily\bfseries\selectfont, coltitle=black,
    fontupper=\sffamily, interior style={left color=tcol_CNT1!80,right
    color=tcol_CNT2!80}, frame style={left color=tcol_CNT1!60!black,right
    color=tcol_CNT2!60!black}, attach boxed title to top center={yshift=10pt},
    boxed title style={frame hidden, interior style={left color=tcol_CNT1,right
    color=tcol_CNT2}, frame style={left color=tcol_CNT1!60!black,right
    color=tcol_CNT2!60!black}, height=24pt,bean arc,drop fuzzy shadow },
    top=2mm,bottom=2mm,left=2mm,right=2mm, before skip=20mm,after skip=20mm,
    drop fuzzy shadow,breakable]
%
\makeatletter
\@starttoc{toc}
\makeatother
\end{tcolorbox}


\section{The Fate of the Earth, and other Exoplanets}
This project will use MESA to simulate the evolution of a Sun-analog star with a
mass of $1\,M_{\odot}$ and metallicity of $Z=0.02$. MESA is the most widely used
open source stellar evolution code. You should first start by following the
\href{https://docs.mesastar.org/en/release-r23.05.1/quickstart.html#installation}{installation
instructions for MESA}. As with all open source codes, it is a good idea to read
all the documentation provided in the `Quickstart' if there is one. Once you've
acquainted yourself with the documentation, start working to build your Sun
analog by following along with the
\href{https://docs.mesastar.org/en/release-r23.05.1/test_suite/1M_pre_ms_to_wd.html#m-pre-ms-to-wd}{test
case for a $1\,M_{\odot}$ star}. What are the conditions for habitability? At
different phases of the Sun's evolution, how do these conditions change? 

Once you have a working stellar evolution model, make a change to the stellar
properties in the way you find most interesting. You'll need to justify the
change you make, so choose carefully! $:-)$ Where would the habitability zone be
for your new stellar model? And how does the star's evolution affect that
habitability? Compare your stars to the hosts of known exoplanets. Why are they
different? Are the differences due to physical reasons or due to our search
strategies for exoplanets? What search strategies are best for finding
`habitable planets'?



\section{Testing Stellar Models with Wide Binaries}
Gaia is a telescope launched by the European Space Agency which is measuring the
positions and velocities of every star in the Milky Way that is brighter than
$G=20$ by observing the entire sky many times over its 10-year observation
duration. This survey design creates a myriad of \emph{extremely rich} datasets
because the time-resolved data allows for 6-dimensional (3 dimensions in
position and 3 dimensions in velocity) measurements!

One early study used Gaia data to construct a catalog of roughly a million wide
binary stars. Multiple projects can be done with this dataset! Since multiple
projects will be taken from the same dataset, students who pick the projects
below will be able to work together on aspects of the projects that are shared.
Any project that is done in this category should begin with a thorough reading
of the associated
\href{https://ui.adsabs.harvard.edu/abs/2021MNRAS.506.2269E/abstract}{paper} by
El-Badry et al. 2022 which links to the data and details how the population was
selected. Each person/team should reselect the data then perform an analysis
chosen from the list below. You should be able to summarize the selection
process and describe why this selection ensures that the binary systems are wide
enough that they are unlikely to have interacted. This project has tool overlap
with the `Measuring Star Cluster Ages' project so you should plan to share
knowledge and expertise with your colleagues working on that project. 

\subsection{Calibrating the Initial to Final Mass Relation for White Dwarfs}
 The initial to final mass relation is a way to connect stars at ZAMS to the WDs
 they produce and is used in a variety of astronomical contexts. The wide binary
 catalog contains several binaries containing two WDs. These are very
 interesting sources for calibration since time since the formation of each WD
 can be calculated from models and the ZAMS formation time of each WD progenitor
 should be the same since they are in a binary system. This means we can solve
 for the time from ZAMS formation to WD formation for each star.
 
 This project will calculate the WD formation times for the double WD binaries
 in the Gaia catalog and then compare those findings to stellar isochrones from
 the \href{https://waps.cfa.harvard.edu/MIST/}{MESA Isochrones and Stellar
 Tracks (MIST)} catalogs to obtain an \emph{empirical} initial to final mass
 relation built by each of the WDs in the binary. You should then compare this
 empirical initial to final mass relation to others reported in the literature.
 How do they compare?

\subsection{Testing Stellar Age Measurements from Isochrones}
Isochrones are built from theoretical models for stellar evolution and show
snapshots of a single-age stellar population at different ages. This means that
they are unlikely to produce age estimates which are $100\%$ accurate. This
project will use binaries containing a white dwarf (WD) and a main sequence (MS)
star. Since WDs have very clearly defined structure, you should first determine
how long ago they formed based on their temperature and brightness by comparing
to cooling models. You should then use the mass of the WD and compare to an
initial to final mass relation to determine the ZAMS mass and total age of the
binary system. You should then obtain the age of the (MS) by comparing it to
isochrones which match the metallicity of the star. How does the age of the MS
star compare to the age of obtained for the WD? Repeat this process for as many
WD - MS binaries as you can to look for trends in the age measurements across MS
mass and metallicity. 

\section{Analyzing Periodic Signals with Lightkurve}
Photometry is the easiest form of astronomical measurement. Time-resolved
photometric data can be used to study populations of stars in great detail! In
order to use photometric data, though, software tools are necessary for
analysis. The following projects will use the Python package:
\href{https://docs.lightkurve.org/index.html}{lightkurve} to carry out analysis
of stars based solely on photometric time-series data. For each project, you
should follow the tutorials which show how to access Kepler and TESS data and to
produce observed light curves of each source you study. 

\subsection{Measuring the Rotation Rates of Stars}
This project will follow
\href{https://docs.lightkurve.org/tutorials/index.html#rotation-rates-and-periodic-signals}{lightkurve
tutorials} which show how to analyze photometric data to measure the rotation
rates of stars and then subtract that rotation period out of the data. There
will often be residual periodic signals; you should search for other sources of
periodic signals beyond stellar rotation and be able to describe each of them. 

Stellar rotation rates are also used to produce age estimates for stellar
populations. You should search the literature to study how age estimates work in
practice with stellar rotation.Once you finish the tutorials which give examples
for a couple of stars, you should extend your skills to measure the rotation
rates of other Kepler stars. 


\subsection{Estimating Stellar Properties with Asteroseismology}
This project will follow a
\href{https://docs.lightkurve.org/tutorials/3-science-examples/asteroseismology-estimating-mass-and-radius.html}{lightkurve
tutorial} which shows how to measure the mass and radius of a star through
asteroseismic photometric measruements. Along with the tutorial, you should do
research on asteroseismology to investigate how it is used to measure stellar
properties that are difficult to obtain otherwise. The lightkurve tutorial uses
a single star observed by Kepler, a space telescope designed to find exoplanets.
Once you've completed the tutorial, you should expand your analysis to other
stars observed by Kepler. By the end of this project you should be able to
describe the key principles of asteroseismology, demonstrate its use for
measuring stellar properties, and explain why it is not a tool that can be used
for stars at all stages of their evolution.  


\section{Measuring Star Cluster Ages}
Star clusters are excellent tools for studying stellar populations because they
contain stars that are all born with (roughly) the same age, composition, and
distance. As a result, the brightness and temperature of the stars in a cluster
can be used to provide an age estimate of the population if the stars in a given
cluster are compared to a stellar isochrone. 

For this project, you should pick a star cluster and obtain data for its
constituent stars. Gaia is a good place to start looking for data! Once you have
the data, you should make an absolute color-magnitude diagram of the population.
You should then use the \href{https://waps.cfa.harvard.edu/MIST/}{MESA
Isochrones and Stellar Tracks (MIST)} catalogs to calculate the age of the
cluster. Based on the age of the cluster and the number of stars shining at
present, how many stars do you expect to have already evolved through their
lives to produce a stellar remnant? This project has tool overlap with the
`Testing Stellar Models with Wide Binaries' projects so you should plan to share
knowledge and expertise with your colleagues working on those projects. 

\section{Designing Surveys for the Stars}
Stellar evolution timescales rarely overlap with human timescales. This means
that in order to test stellar evolution models, we need surveys that can observe
stars and stellar populations across all phases of their lives. Unfortunately,
funding for astronomy is limited and we must choose carefully which surveys we
decide to support as a community.  

For this project, you should explore different strategies for observing stars at
evolutionary phases across the H-R Diagram. You should identify current surveys
which are already planned or taking data now. Are there any phases of evolution
that are not well covered by current or planned surveys? If you were given a
free pass for one ground-based or space-based survey, how would you invest. You
should be able to provide strong reasoning for your choice including why that
investment would help other fields of astoronomy. 


\section{Supernova Natal Kicks}
Supernovae are some of the most energetic events in the Universe. Most
supernovae occur at the end of a massive star's life once fusion in the core has
ceased. For this project you should first search the literature to investigate
different theoretical models for core-collapse and electron capture supernovae.
How do these models influence the natal kicks imparted to the newly formed
compact objects? A good starting place to search for models that are widely used
in the literature is the documentation for the
\href{https://cosmic-popsynth.github.io/}{COSMIC} population synthesis code
which incorporates different assumptions for supernovae. 

Next, you should investigate the literature for astronomical datasets that can
be used to constrain natal kicks. Some starting ideas are neutron stars in
supernova remnants or pulsars in globular clusters. Pick a dataset to focus on
to try to determine whether the compact objects are likely to have received weak
or strong natal kicks. You should be able to support your arguments with
quantitative evidence that shows how supernova model assumptions compare to the
data. 

\section{The Black Hole Mass Function}
Black holes with masses below $\sim100\,M_{\odot}$ are predominantly expected to
form from stars. This project will investigate how different assumptions for
stellar evolution affects the masses of the black holes that can be produced
from stars by using the population synthesis code
\href{https://cosmic-popsynth.github.io/}{COSMIC} to simulate single stars with
different ZAMS masses and metallicities. COSMIC includes several models for
compact object formation but actually doesn't treat black holes different from
neutron stars; your first task to study how COSMIC simulations compact object
formation and articulate why black holes and neutron stars are handled the same
way. 

Next, you should investigate the boundaries (upper and lower) of the black hole
mass spectrum and how those boundaries change with different assumptions in the
code. How does the mass spectrum within the boundaries change? Compare your
simulations to black holes discovered and characterized through microlensing,
radial velocity and astrometry, X-rays, and gravitational waves. What can you
conclude about the origin of black holes based on this comparison? As detectors
increase their sensitivity to stellar populations that are far away, do you
expect the majority of the black holes to have larger masses or smaller masses?



\end{document}